\begin{abstract}
Timing channels present a significant security threat in modern computing,
allowing unauthorized information flow through observable variations in
execution time. Research has been done on creating secure microkernels designed
to mitigate these channels from the ground up \cite{ge2019, karlsson2025}, while
other research focuses on closing timing channels at the user
level \cite{cauligi2019}. It remains an open challenge to gather these findings
for a more general and mainstream OS, and this thesis aims to investigate the
feasibility of integrating time protection mechanisms into the widely used
FreeRTOS kernel.

FreeRTOS provides a default priority-based preemptive scheduler, which is
inherently vulnerable to timing channels. To address this, the project
implements a \textbf{domain-level} scheduler that partitions execution into
fixed time slices, ensuring that the execution of tasks in one domain does
not affect the timing observed by another. Key to this implementation is the
temporal fence (fence.t) instruction during domain switches to clear shared
microarchitectural states \cite{wistoff2023, wistoff2025}. Further,
constant-time domain creation mechanisms are introduced, allowing dynamic
resource reconfiguration without introducing new timing channels.

The implementation was evaluated using QEMU system emulation. The results
demonstrate that the modified kernel achieves constant-time domain scheduling
with negligible variance with regard to instruction counting. While evaluation
on true hardware remains undone, this work provides a foundation towards
bringing timing protection to mainstream embedded systems.
\end{abstract}
